\section*{ID6763: F-Number = Bessel Behavior = Directional Aperture Geometry}
\paragraph{Metadata.}
PiID: 3859381023439 | RTEID: RTE:P26262.2309:T5.5730 | UUID: 8716abc6-26b2-5da2-bfa9-72afa0d78a3e

\paragraph{Canonical Statement.}
\textbackslash{}[ I(\textbackslash{}theta) \textbackslash{}propto \textbackslash{}sum\_\{m=0\}\textasciicircum{}\{N-1\} \textbackslash{}left| J\_m(k R) e\textasciicircum{}\{i m \textbackslash{}theta\} \textbackslash{}right|\textasciicircum{}2 \textbackslash{}]

\paragraph{Canonical Equation.}
\[
I(\theta) \propto \sum_{m=0}^{N-1} \left| J_m(k R) e^{i m \theta} \right|^2
\]

\paragraph{Dictionary.}
F-Number = Bessel Behavior = Directional Aperture Geometry — I(θ) ∝ ∑\_m=0\textasciicircum{}N-1 | J\_m(k R) e\textasciicircum{}i m θ |\textasciicircum{}2

\paragraph{Encyclopedia.}
F-Number = Bessel Behavior = Directional Aperture Geometry is a series / summation, written I(θ) ∝ ∑\_m=0\textasciicircum{}N-1 | J\_m(k R) e\textasciicircum{}i m θ |\textasciicircum{}2. It is expressed as a sum or series over modes or indices, superposing discrete contributions. It is built from θ, I, m, N, J\_m, k. Within the Trawin Zero Point registry it serves as one element of the framework's mathematical narrative.
